% This chapter was modified on 1/12/05.
%\setcounter{chapter}{7}
\chapter{Hukum Bilangan Besar}\label{chp 8} 

\section[Peubah Acak Diskret]{Hukum Bilangan Besar untuk Peubah Acak Diskret}
\label{sec 8.1}
Kita kini siap membuktikan teorema fundamental pertama kita dalam probabilitas.  
Kita telah melihat bahwa cara intuitif untuk memandang peluang suatu hasil tertentu adalah sebagai 
frekuensi kemunculan hasil tersebut dalam jangka panjang, ketika percobaan diulang sangat 
banyak kali.  Kita juga telah mendefinisikan peluang secara matematis sebagai nilai suatu fungsi 
distribusi bagi peubah acak yang merepresentasikan percobaan.   Hukum Bilangan Besar, yaitu
teorema yang dibuktikan dalam model matematis probabilitas, menunjukkan bahwa model ini konsisten 
dengan interpretasi frekuensi atas peluang.  Teorema ini kadang-kadang disebut \emx {hukum
rata-rata.}  Untuk mengetahui apa yang akan terjadi jika hukum ini tidak benar, lihat artikel karya 
Robert M. Coates.\index{COATES, R. M.}\footnote{R.~M.~Coates, ``The Law," \emx {The World of
Mathematics,} ed. James R. Newman (New York: Simon and Schuster, 1956.}

\subsection*{Ketaksamaan Chebyshev}
Untuk membahas Hukum Bilangan Besar, pertama-tama kita memerlukan sebuah ketaksamaan penting
yang disebut \emx {Ketaksamaan Chebyshev.}

\begin{theorem}{\bf (Ketaksamaan Chebyshev)}\index{Ketaksamaan Chebyshev}
Misalkan $X$ adalah peubah acak diskret dengan nilai harapan $\mu = E(X)$, dan misalkan $\epsilon
> 0$ sebarang bilangan real positif.  Maka
$$
P(|X - \mu| \geq \epsilon) \leq \frac {V(X)}{\epsilon^2}\ .
$$
\proof
Misalkan $m(x)$ menyatakan fungsi distribusi $X$.  Maka peluang bahwa $X$
berselisih dari $\mu$ sekurang-kurangnya sebesar $\epsilon$ diberikan oleh
$$P(|X - \mu| \geq \epsilon) = \sum_{|x - \mu| \geq \epsilon} m(x)\ .$$
Kita mengetahui bahwa 
$$V(X) = \sum_x (x - \mu)^2 m(x)\ ,$$
dan nilai ini jelas sekurang-kurangnya sebesar
$$\sum_{|x - \mu| \geq \epsilon} (x - \mu)^2 m(x)\ ,$$
karena semua suku penjumlahan bernilai positif dan kita telah membatasi rentang penjumlahan pada
jumlah kedua.  Namun, jumlah terakhir ini sekurang-kurangnya sebesar
\begin{eqnarray*}
\sum_{|x - \mu| \geq \epsilon} \epsilon^2 m(x) &=& 
\epsilon^2 \sum_{|x - \mu| \geq \epsilon} m(x) \\
&=& \epsilon^2 P(|X - \mu| \geq \epsilon)\ .\\
\end{eqnarray*}
Jadi,
$$ P(|X - \mu| \geq \epsilon) \leq \frac {V(X)}{\epsilon^2}\ .
$$
\end{theorem}
Perhatikan bahwa $X$ dalam teorema di atas dapat berupa sebarang peubah acak diskret, dan
$\epsilon$ dapat berupa sebarang bilangan positif.

\begin{example}
Misalkan $X$ sebarang peubah acak dengan $E(X) = \mu$ dan $V(X) = \sigma^2$.  Maka,
jika $\epsilon = k\sigma$, Ketaksamaan Chebyshev menyatakan bahwa
$$
P(|X - \mu| \geq k\sigma) \leq \frac {\sigma^2}{k^2\sigma^2} = \frac 1{k^2}\ .
$$
Jadi, untuk sebarang peubah acak, peluang penyimpangan dari rata-rata yang besarnya
lebih dari~$k$ simpangan baku adalah ${} \leq 1/k^2$.  Sebagai contoh, jika
$k = 5$, $1/k^2 = .04$.
\end{example}

Ketaksamaan Chebyshev merupakan ketaksamaan terbaik yang mungkin dalam arti bahwa, untuk
sebarang $\epsilon > 0$, kita dapat memberikan contoh peubah acak yang
membuat Ketaksamaan Chebyshev berlaku sebagai kesamaan.  Untuk melihat hal ini, dengan
$\epsilon > 0$ yang diberikan, pilih $X$ dengan distribusi
$$
p_X = \pmatrix{
-\epsilon & +\epsilon \cr
1/2 & 1/2 \cr}\ .
$$
Maka $E(X) = 0$, $V(X) = \epsilon^2$, dan
$$
P(|X - \mu| \geq \epsilon) = \frac {V(X)}{\epsilon^2} = 1\ .
$$

Kita kini siap menyatakan dan membuktikan Hukum Bilangan Besar.

\subsection*{Hukum Bilangan Besar}
\begin{theorem}{\bf (Hukum Bilangan Besar)}\index{Hukum Bilangan Besar}
Misalkan $X_1$,~$X_2$, \dots,~$X_n$ membentuk proses percobaan saling bebas, dengan
nilai harapan berhingga $\mu = E(X_j)$ dan varians berhingga $\sigma^2 =
V(X_j)$.  Misalkan $S_n = X_1 + X_2 +\cdots+ X_n$.  Maka untuk
sebarang $\epsilon > 0$, 
$$ P\left( \left| \frac {S_n}n - \mu \right| \geq \epsilon
\right) \to 0
$$
ketika $n \rightarrow \infty$.
Secara ekuivalen,
$$
P\left( \left| \frac {S_n}n - \mu \right| < \epsilon \right) \to 1
$$
ketika $n \rightarrow \infty$.
\proof
Karena $X_1$,~$X_2$, \dots,~$X_n$ saling bebas dan memiliki distribusi yang sama,
kita dapat menerapkan Teorema~\ref{thm 6.9}.  Kita memperoleh
$$
V(S_n) = n\sigma^2\ ,
$$
dan
$$
V (\frac {S_n}n) = \frac {\sigma^2}n\ .
$$
Kita juga mengetahui bahwa
$$
E (\frac {S_n}n) = \mu\ .
$$
Berdasarkan Ketaksamaan Chebyshev, untuk sebarang $\epsilon > 0$,
$$
P\left( \left| \frac {S_n}n - \mu \right| \geq \epsilon \right) \leq \frac
{\sigma^2}{n\epsilon^2}\ .
$$
Jadi, untuk $\epsilon$ tetap,
$$
P\left( \left| \frac {S_n}n - \mu \right| \geq \epsilon \right) \to 0
$$
ketika $n \rightarrow \infty$, atau secara ekuivalen,
$$
P\left( \left| \frac {S_n}n - \mu \right| < \epsilon \right) \to 1
$$
ketika $n \rightarrow \infty$.
\end{theorem}

\subsection*{Hukum Rata-Rata}
Perhatikan bahwa $S_n/n$ adalah rata-rata hasil-hasil individual, dan Hukum
Bilangan Besar sering disebut ``hukum rata-rata."  Merupakan kenyataan yang mencolok
bahwa kita dapat memulai dengan percobaan acak yang hanya sedikit dapat diprediksi
dan, dengan mengambil rata-rata, memperoleh percobaan yang hasilnya dapat
diprediksi dengan tingkat kepastian tinggi.  Hukum Bilangan Besar, sebagaimana
telah kita nyatakan, sering disebut ``Hukum Lemah Bilangan Besar" untuk
membedakannya dari ``Hukum Kuat Bilangan Besar" yang dijelaskan dalam
Latihan~\ref{exer 8.1.16}.

Pertimbangkan kasus khusus penting berupa percobaan Bernoulli dengan peluang~$p$
untuk sukses.  Misalkan $X_j = 1$ jika hasil ke-$j$ adalah sukses dan~0 jika
gagal.  Maka $S_n = X_1 + X_2 +\cdots+ X_n$ adalah banyaknya sukses dalam $n$
percobaan dan $\mu = E(X_1) = p$.  Hukum Bilangan Besar menyatakan bahwa untuk sebarang
$\epsilon > 0$
$$
P\left( \left| \frac {S_n}n - p \right| < \epsilon \right) \to 1
$$
ketika $n \rightarrow \infty$.  Pernyataan di atas mengatakan bahwa, dalam pengulangan
percobaan Bernoulli yang jumlahnya besar, kita dapat mengharapkan proporsi percobaan yang
menghasilkan kejadian tersebut mendekati $p$.  Ini menunjukkan bahwa model matematis probabilitas kita
sesuai dengan interpretasi frekuensi atas peluang.  

\subsection*{Pelemparan Koin}
Mari kita pertimbangkan kasus khusus pelemparan sebuah koin sebanyak $n$ kali, dengan $S_n$
menyatakan banyaknya kemunculan sisi kepala.  Maka peubah acak $S_n/n$ menyatakan
proporsi kemunculan sisi kepala dan memiliki nilai antara 0~dan~1.  Hukum
Bilangan Besar memprediksi bahwa hasil peubah acak ini, untuk
$n$ yang besar, akan mendekati 1/2.

Pada Gambar~\ref{fig 8.1}, kita telah menggambar distribusi dalam contoh ini untuk nilai
~$n$ yang semakin besar. Kita telah menandai hasil antara .45~dan~.55 dengan titik-titik di puncak
batang.  Kita melihat bahwa ketika
$n$ meningkat, distribusi semakin terkonsentrasi di sekitar~.5 dan persentase
luas total yang termuat dalam interval $(.45,.55)$ semakin
besar, sebagaimana diprediksi oleh Hukum Bilangan Besar.

\putfig{5.0truein}{PSfig8-1}{Distribusi percobaan Bernoulli.}{fig 8.1} 


\subsection*{Pelemparan Dadu}
\begin{example}
Pertimbangkan $n$ lemparan sebuah dadu.  Misalkan $X_j$ adalah hasil lemparan ke-$j$.  Maka
$S_n = X_1 + X_2 +\cdots+ X_n$ adalah jumlah hasil dari $n$ lemparan pertama.  Ini merupakan
proses percobaan saling bebas dengan $E(X_j) = 7/2$.  Jadi, berdasarkan Hukum Bilangan
Besar, untuk sebarang $\epsilon > 0$
$$
P\left( \left| \frac {S_n}n - \frac 72 \right| \geq \epsilon \right) \to 0
$$
ketika $n \rightarrow \infty$.  Pernyataan yang ekuivalen adalah bahwa, untuk sebarang $\epsilon > 0$,
$$
P\left( \left| \frac {S_n}n - \frac 72 \right| < \epsilon \right) \to 1
$$
ketika $n \rightarrow \infty$.
\end{example}

\subsection*{Perbandingan Numerik}
Perlu ditekankan bahwa, walaupun Ketaksamaan Chebyshev membuktikan Hukum Bilangan
Besar, ketaksamaan tersebut sebenarnya sangat kasar untuk peluang-peluang yang terlibat.  Namun,
kekuatannya terletak pada kenyataan bahwa ketaksamaan tersebut berlaku bagi sebarang peubah acak dan memungkinkan
kita membuktikan teorema yang sangat kuat.
\par
Dalam contoh berikut, kita membandingkan taksiran yang diberikan oleh Ketaksamaan Chebyshev dengan
nilai sebenarnya.
\begin{example}
Misalkan $X_1$,~$X_2$, \dots,~$X_n$ membentuk proses percobaan Bernoulli dengan
peluang~.3 untuk sukses dan~.7 untuk kegagalan.  Misalkan $X_j = 1$ jika hasil ke-$j$
adalah sukses dan~0 jika gagal.  Maka, $E(X_j) = .3$ dan $V(X_j) =
(.3)(.7) = .21$.  Jika
$$
A_n = \frac {S_n}n = \frac {X_1 + X_2 +\cdots+ X_n}n
$$
adalah \emx {rata-rata} dari $X_i$, maka $E(A_n) = .3$ dan $V(A_n) =
V(S_n)/n^2 = .21/n$.  Ketaksamaan Chebyshev menyatakan bahwa, jika misalnya
$\epsilon = .1$,
$$
P(|A_n - .3| \geq .1) \leq \frac {.21}{n(.1)^2} = \frac {21}n\ .
$$
Jadi, jika $n = 100$,
$$
P(|A_{100} - .3| \geq .1) \leq .21\ ,
$$
atau jika $n = 1000$,
$$
P(|A_{1000} - .3| \geq .1) \leq .021\ .
$$
Pernyataan ini dapat ditulis ulang sebagai
 \begin{eqnarray*}
P(.2 < A_{100} < .4) &\geq& .79\ , \\
P(.2 < A_{1000} < .4) &\geq& .979\ .
\end{eqnarray*} 
Nilai-nilai ini perlu dibandingkan dengan nilai sebenarnya, yaitu (hingga enam tempat
desimal)
\begin{eqnarray*}
P(.2 < A_{100} < .4) &\approx& .962549 \\
P(.2 < A_{1000} < .4) &\approx& 1\ .\\
\end{eqnarray*}
Program {\bf Law}\index{Law (program)} dapat digunakan untuk melakukan perhitungan di atas secara
sistematis.
\end{example}


\subsection*{Catatan Historis}
Hukum Bilangan Besar pertama kali dibuktikan oleh matematikawan Swiss, James
Bernoulli\index{BERNOULLI, J.|(}, dalam bagian keempat karyanya, \emx {Ars Conjectandi}, yang diterbitkan
secara anumerta pada~1713.\footnote{J. Bernoulli, \emx {The Art of Conjecturing
IV,} trans.~Bing Sung, Technical Report No.~2, Dept.\ of Statistics, Harvard
Univ., 1966}  Seperti yang sering terjadi pada pembuktian pertama, pembuktian Bernoulli jauh
lebih sulit daripada pembuktian yang telah kita sajikan menggunakan Ketaksamaan Chebyshev. 
Chebyshev mengembangkan ketaksamaannya untuk membuktikan bentuk umum Hukum Bilangan
Besar (lihat Latihan~\ref{exer 8.1.13}).  Ketaksamaan itu sendiri muncul
jauh lebih awal dalam karya Bienaym\'e.\index{BIENAYM\'E, I.} Dalam pembahasannya tentang sejarah ketaksamaan ini,
Maistrov\index{MAISTROV, L.} menyatakan bahwa ketaksamaan tersebut lama disebut Ketaksamaan
Bienaym\'e-Chebyshev.\footnote{L. E. Maistrov, \emx {Probability Theory: A Historical
Approach,} trans.\ and ed.~Samual Kotz, (New York: Academic Press, 1974),
p.~202}

Dalam \emx {Ars Conjectandi}, Bernoulli menyajikan kepada pembacanya pembahasan panjang
tentang makna teoremanya disertai banyak contoh.  Dalam notasi modern, ia membahas
sebuah kejadian yang terjadi dengan peluang~$p$, tetapi ia tidak mengetahui $p$.  Ia ingin
menaksir $p$ menggunakan $\bar{p}$, yakni proporsi pengulangan yang menghasilkan kejadian tersebut,
dari sejumlah pengulangan percobaan.  Ia membahas secara terperinci
masalah menaksir, dengan metode ini, proporsi bola putih dalam sebuah guci
yang memuat bola putih dan bola hitam dalam jumlah yang tidak diketahui.  Ia melakukannya dengan
mengambil serangkaian bola dari guci, mengembalikan bola yang diambil setelah setiap
pengambilan, dan menaksir proporsi bola putih yang tidak diketahui dalam guci dengan
proporsi bola putih di antara bola-bola yang diambil.  Ia menunjukkan bahwa, dengan memilih $n$
cukup besar, ia dapat memperoleh tingkat ketepatan dan keandalan apa pun yang diinginkan untuk
taksiran tersebut.  Ia juga memberikan pembahasan yang menarik tentang penerapan
teoremanya untuk menaksir peluang meninggal akibat penyakit tertentu,
terjadinya berbagai jenis cuaca, dan sebagainya.

Ketika membicarakan banyaknya percobaan yang diperlukan untuk membuat suatu penilaian, Bernoulli
mengamati bahwa bahkan tanpa pelatihan formal orang memahami dua hal: satu atau dua percobaan
tidak memadai, dan risiko taksiran meleset dari sasaran berkurang ketika jumlah pengamatan
diperbesar.\footnote{Bernoulli, op.\ cit., p.~38. Ringkasan editorial; ungkapan terjemahan
modern tidak diterjemahkan atau direproduksi.}

\noindent Namun, ia juga merumuskan pertanyaan yang harus dijawab: ketika banyaknya
pengamatan bertambah, apakah peluang bahwa frekuensi teramati mendekati rasio sebenarnya dapat
melampaui setiap taraf kepastian yang ditentukan sebelumnya, ataukah terdapat batas atas yang
tidak dapat dilampaui?\footnote{ibid., p.~39. Ringkasan editorial; ungkapan terjemahan
modern tidak diterjemahkan atau direproduksi.}

\noindent Bernoulli menekankan bahwa ia telah memikirkan masalah tersebut selama dua
puluh tahun dan menilai kebaruan, kegunaan, serta kesulitannya cukup besar untuk menandingi
seluruh bagian lain dalam risalahnya.\footnote{ibid., p.~42. Ringkasan editorial; ungkapan
terjemahan modern tidak diterjemahkan atau direproduksi.}

\noindent Pada akhir pembuktiannya, Bernoulli menghubungkan pengamatan tanpa batas dengan
munculnya rasio tetap dan pola pergiliran yang stabil: bahkan kejadian acak akan tampak
mengikuti suatu keniscayaan. Ia lalu membandingkan gagasan itu dengan ajaran Plato tentang
kembalinya segala sesuatu secara universal setelah rentang zaman yang amat panjang.\footnote{ibid.,
pp.~65--66. Ringkasan editorial; ungkapan terjemahan modern tidak diterjemahkan atau
direproduksi.}\index{BERNOULLI, J.|)}

\exercises
\begin{LJSItem}

\i\label{exer 8.1.1} Sebuah koin seimbang dilempar 100 kali.  Nilai harapan banyaknya kemunculan
sisi kepala adalah~50, dan simpangan baku bagi banyaknya kemunculan sisi kepala adalah $(100 \cdot
1/2 \cdot 1/2)^{1/2} = 5$.  Apakah yang dinyatakan Ketaksamaan Chebyshev tentang
peluang bahwa banyaknya kemunculan sisi kepala menyimpang dari nilai harapan
50 sebesar tiga simpangan baku atau lebih (yakni, sekurang-kurangnya 15)?

\i\label{exer 8.1.100} Tulis sebuah program yang menggunakan fungsi 
$\mbox {binomial}(n,p,x)$ untuk menghitung secara tepat peluang yang Anda taksir dalam
Latihan~\ref{exer 8.1.1}.  Bandingkan kedua hasilnya.

\i\label{exer 8.1.101} Tulis sebuah program untuk melempar koin 10{,}000 kali.  Misalkan $S_n$ adalah
banyaknya kemunculan sisi kepala dalam $n$ lemparan pertama.  Mintalah program Anda mencetak, setelah setiap
1000 lemparan, $S_n - n/2$.  Berdasarkan simulasi ini, apakah tepat
mengatakan bahwa Anda dapat mengharapkan sisi kepala muncul kira-kira separuh waktu ketika melempar koin
berkali-kali?

\i\label{exer 8.1.102} Taruhan 1 dolar pada permainan craps memiliki nilai harapan keuntungan $-.0141$.  Apa
yang dikatakan Hukum Bilangan Besar tentang keuntungan Anda jika Anda memasang banyak
taruhan 1 dolar di meja craps?  Apakah hukum tersebut menjamin bahwa kerugian Anda akan
kecil?  Apakah hukum tersebut menjamin bahwa jika $n$ sangat besar Anda akan mengalami kerugian?

\i\label{exer 8.1.103} Misalkan $X$ adalah peubah acak dengan $E(X) =0$ dan $V(X) = 1$.  Nilai
bilangan bulat~$k$ berapakah yang menjamin bahwa $P(|X| \geq k) \leq .01$?

\i\label{exer 8.1.6} Misalkan $S_n$ adalah banyaknya sukses dalam $n$ percobaan
Bernoulli dengan peluang~$p$ untuk sukses pada setiap percobaan.  Dengan menggunakan Ketaksamaan
Chebyshev, tunjukkan bahwa untuk sebarang $\epsilon > 0$
$$
P\left( \left| \frac {S_n}n - p \right| \geq \epsilon \right) \leq \frac {p(1 -
p)}{n\epsilon^2}\ .
$$

\i\label{exer 8.1.7} Tentukan nilai maksimum yang mungkin bagi $p(1 - p)$ jika $0 < p
< 1$.  Dengan menggunakan hasil ini dan Latihan~\ref{exer 8.1.6}, tunjukkan bahwa taksiran
$$
P\left( \left| \frac {S_n}n - p \right| \geq \epsilon \right) \leq \frac
1{4n\epsilon^2}
$$
berlaku untuk sebarang $p$.

\i\label{exer 8.1.104} Sebuah koin seimbang dilempar berkali-kali.  Apakah Hukum Bilangan
Besar menjamin bahwa, jika $n$ cukup besar, dengan $\mbox {peluang} >
.99$ banyaknya kemunculan sisi kepala tidak akan menyimpang dari $n/2$ lebih dari
100?

\i\label{exer 8.1.105} Dalam Latihan~\ref{sec 6.2}.\ref{exer 6.2.16}, Anda menunjukkan bahwa, untuk
masalah penitipan topi, banyaknya orang $S_n$ yang memperoleh kembali topi mereka sendiri memiliki $E(S_n) =
V(S_n) = 1$.  Dengan menggunakan Ketaksamaan Chebyshev, tunjukkan bahwa $P(S_n \geq 11) \leq .01$
untuk sebarang $n \geq 11$.

\i\label{exer 8.1.106} Misalkan $X$ sebarang peubah acak yang mengambil nilai 0,~1,~2,
\dots,~$n$ dan memiliki $E(X) = V(X) = 1$.  Tunjukkan bahwa, untuk sebarang bilangan bulat positif $k$,
$$
P(X \geq k + 1) \leq \frac 1{k^2}\ .
$$

\i\label{exer 8.1.107} Kita mempunyai dua koin: satu koin seimbang dan satu lagi menghasilkan
sisi kepala dengan peluang 3/4.  Salah satu dari kedua koin dipilih secara acak,
dan koin tersebut dilempar $n$ kali.  Misalkan $S_n$ adalah banyaknya kemunculan sisi kepala
dalam $n$ lemparan ini.  Apakah Hukum Bilangan Besar memungkinkan kita memprediksi
proporsi kemunculan sisi kepala dalam jangka panjang?  Setelah mengamati
banyak lemparan, dapatkah kita mengetahui koin mana yang dipilih?  Berapa banyak lemparan
yang cukup untuk membuat kita yakin 95~persen?

\i\label{exer 8.1.13} (Chebyshev\index{CHEBYSHEV, P. L.}\footnote{P. L. Chebyshev, ``On Mean
Values," \emx {J.\ Math.\ Pure.\ Appl.,} vol.~12 (1867), pp.~177--184.})  Misalkan
$X_1$,~$X_2$, \dots,~$X_n$ adalah peubah acak saling bebas yang mungkin memiliki
distribusi berbeda, dan misalkan $S_n$ adalah jumlahnya.  Misalkan $m_k = E(X_k)$,
$\sigma_k^2 = V(X_k)$, dan $M_n = m_1 + m_2 +\cdots+ m_n$.  Andaikan
$\sigma_k^2 < R$ untuk semua~$k$.  Buktikan bahwa, untuk sebarang $\epsilon > 0$,
$$
P\left( \left| \frac {S_n}n - \frac {M_n}n \right| < \epsilon \right) \to 1
$$
ketika $n \rightarrow \infty$.

\i\label{exer 8.1.108}  Sebuah koin seimbang dilempar berulang kali.  Sebelum setiap lemparan, Anda boleh
memutuskan apakah akan bertaruh pada hasilnya.  Dapatkah Anda menjelaskan sistem taruhan dengan
tak terhingga banyak taruhan yang memungkinkan Anda, dalam jangka panjang, memenangkan lebih dari separuh taruhan? 
(Perhatikan bahwa kita tidak mengizinkan sistem taruhan yang menyuruh Anda bertaruh sampai unggul, lalu
berhenti.) Tulis sebuah program komputer yang menerapkan sistem taruhan ini.  Seperti dinyatakan di atas,
program Anda harus memutuskan apakah akan bertaruh pada suatu hasil tertentu sebelum hasil tersebut ditentukan.
Sebagai contoh, Anda dapat memilih hanya hasil yang muncul setelah
sisi ekor muncul tiga kali berturut-turut.  Uji apakah Anda dapat memperoleh lebih dari 50\% sisi kepala
dengan ``sistem" Anda.

\istar\label{exer 8.1.109} Buktikan analog berikut dari Ketaksamaan Chebyshev:
$$
P(|X - E(X)| \geq \epsilon) \leq \frac 1\epsilon E(|X - E(X)|)\ .
$$

\istar\label{exer 8.1.16} Kita telah membuktikan teorema yang sering disebut ``Hukum Lemah
Bilangan Besar."  Intuisi kebanyakan orang dan simulasi komputer kita menunjukkan
bahwa, jika kita melempar koin berulang kali, proporsi kemunculan sisi kepala benar-benar
akan mendekati 1/2; yaitu, jika $S_n$ adalah banyaknya kemunculan sisi kepala dalam $n$ lemparan, maka
kita akan memperoleh
$$
A_n = \frac {S_n}n \to \frac 12
$$
ketika $n \to \infty$.  Tentu saja, kita tidak dapat memastikan hal ini karena kita tidak mampu
melempar koin tak terhingga kali dan, seandainya mampu, koin tersebut dapat
selalu menampilkan sisi kepala.  Namun, ``Hukum Kuat Bilangan Besar,"\index{Hukum Kuat 
Bilangan\\ Besar} yang dibuktikan dalam mata kuliah lebih lanjut menyatakan bahwa
$$
P\left( \frac {S_n}n \to \frac 12 \right) = 1\ .
$$
Jelaskan sebuah ruang sampel $\Omega$ yang memungkinkan kita membicarakan
kejadian
$$
E = \left\{\, \omega : \frac {S_n}n \to \frac 12\, \right\}\ .
$$
Bisakah kita menetapkan ukuran berpeluang sama pada ruang ini?  
\choice{}{(Lihat
Contoh~\ref{exam 2.2.12}.)} 

\istar\label{exer 8.1.16.5} 
Dalam latihan ini, kita akan membangun contoh barisan peubah acak
yang memenuhi Hukum Lemah Bilangan Besar, tetapi tidak memenuhi Hukum
Kuat. Distribusi $X_i$ harus bergantung pada $i$, karena
jika tidak, kedua hukum tersebut akan terpenuhi.  (Masalah ini disampaikan kepada kami
oleh David Maslen.)
\vskip .1in
Misalkan kita memiliki barisan tak hingga kejadian yang saling bebas $A_1,
A_2, \ldots$. Misalkan $a_i = P(A_i)$, dan misalkan $r$ adalah bilangan bulat positif.

\begin{enumerate}
\item Tentukan ekspresi untuk peluang bahwa tidak satu pun $A_i$ dengan
$i>r$ terjadi.

\item Gunakan fakta bahwa $x-1 \leq e^{-x}$ untuk menunjukkan bahwa
$$
P(\mbox{Tiada\ $A_i$\ dengan\ $i > r$\ terjadi}) \leq e^{-\sum_{i=r}^{\infty}
a_i}
$$
\item (Lemma Borel-Cantelli pertama)  Buktikan bahwa jika $\sum_{i=1}^{\infty} a_i$ divergen, maka
$$
P(\mbox{Tak\ terhingga\ banyaknya $A_i$\ terjadi}) = 1.
$$
\vskip .1in
\noindent
Sekarang, misalkan $X_i$ adalah barisan peubah acak yang saling bebas sedemikian
sehingga untuk setiap bilangan bulat positif $i \geq 2$,
$$
P(X_i = i) = \frac{1}{2i\log i}, \quad P(X_i = -i) =
\frac{1}{2i\log i}, \quad P(X_i =0) = 1 - \frac{1}{i \log i}.
$$
Ketika $i=1$, kita tetapkan $X_i=0$ dengan peluang $1$.  Seperti biasa, kita tetapkan $S_n = X_1
+ \cdots + X_n$.  Perhatikan bahwa nilai harapan setiap $X_i$ adalah $0$.
\item Tentukan varians $S_n$.
\item Tunjukkan bahwa barisan $\langle X_i \rangle$ memenuhi Hukum Lemah
Bilangan Besar, yakni buktikan bahwa untuk sebarang $\epsilon > 0$
$$
P\biggl(\biggl|{\frac{S_n}{n}}\biggr| \geq \epsilon\biggr) \rightarrow 0\ ,
$$
ketika $n$ menuju tak hingga.
\vskip .1in
\noindent
Sekarang kita menunjukkan bahwa $\{ X_i \}$ tidak memenuhi Hukum Kuat
Bilangan Besar.  Andaikan $S_n / n \rightarrow 0$. Maka, karena
$$
\frac{X_n}{n} = \frac{S_n}{n} - \frac{n-1}{n} \frac{S_{n-1}}{n-1}\ ,
$$
kita mengetahui bahwa $X_n / n \rightarrow 0$.  Dari definisi limit, kita
menyimpulkan bahwa ketaksamaan $|X_i| \geq \frac{1}{2} i$ hanya dapat
berlaku untuk berhingga banyak $i$.
\item Misalkan $A_i$ adalah kejadian $|X_i| \geq \frac{1}{2} i$. Tentukan
$P(A_i)$. Tunjukkan bahwa $\sum_{i=1}^{\infty} P(A_i)$ divergen (gunakan 
Uji Integral).
\item Buktikan bahwa $A_i$ terjadi untuk tak terhingga banyak $i$.
\item Buktikan bahwa 
$$
P\biggl(\frac{S_n}{n} \rightarrow 0\biggr) = 0,
$$
dan karena itu Hukum Kuat Bilangan Besar tidak berlaku bagi barisan
$\{ X_i \}$.

\end{enumerate}
\istar\label{exer 8.1.110} Mari kita lempar sebuah koin bias yang menghasilkan sisi kepala dengan peluang~$p$
dan andaikan Hukum Kuat Bilangan Besar berlaku sebagaimana dijelaskan
dalam Latihan~\ref{exer 8.1.16}.  Maka, dengan peluang~1, $$
\frac {S_n}n \to p
$$
ketika $n \to \infty$.  Jika $f(x)$ adalah fungsi kontinu pada interval satuan,
maka kita juga memperoleh
$$
f\left( \frac {S_n}n \right) \to f(p)\ .
$$

Akhirnya, kita dapat berharap bahwa
$$
E\left(f\left( \frac {S_n}n \right)\right) \to E(f(p)) = f(p)\ .
$$
Tunjukkan bahwa, jika semua ini benar, dan memang demikian, kita telah membuktikan bahwa
setiap fungsi kontinu pada interval satuan merupakan limit fungsi-fungsi
polinomial.  Ini adalah sketsa pembuktian probabilistik bagi sebuah teorema penting 
dalam matematika yang disebut \emx {teorema aproksimasi Weierstrass.}\index{Teorema
Aproksimasi Weierstrass}
\end{LJSItem}


\choice{}{\section[Peubah Acak Kontinu]{Hukum Bilangan Besar untuk Peubah Acak
Kontinu}
\label{sec 8.2}
Dalam bagian sebelumnya, kita membahas dengan cukup terperinci Hukum Bilangan Besar
untuk distribusi peluang diskret.  Hukum ini memiliki analog alami untuk
distribusi peluang kontinu, yang kita bahas secara agak lebih ringkas
di sini.

\subsection*{\bf Ketaksamaan Chebyshev} \hfill\break\index{Ketaksamaan Chebyshev}
Seperti dalam kasus diskret, kita memulai pembahasan dengan Ketaksamaan
Chebyshev.

\begin{theorem}{\bf (Ketaksamaan Chebyshev)} \index{Ketaksamaan Chebyshev}
Misalkan $X$ adalah peubah acak kontinu dengan fungsi densitas $f(x)$.  Andaikan $X$ memiliki 
nilai harapan berhingga $\mu = E(X)$ dan varians berhingga $\sigma^2 = V(X)$.  Maka untuk sebarang
bilangan positif $\epsilon > 0$ berlaku
$$
P(|X - \mu| \geq \epsilon) \leq \frac {\sigma^2}{\epsilon^2}\ .
$$
\end{theorem}

Pembuktiannya sepenuhnya analog dengan pembuktian dalam kasus diskret sehingga tidak kita sertakan.
\par
Perhatikan bahwa teorema ini tidak memberikan informasi jika $\sigma^2 = V(X)$ tak hingga.


\begin{example}
Misalkan $X$ sebarang peubah acak kontinu dengan $E(X) = \mu$ dan $V(X) =
\sigma^2$.  Maka, jika $\epsilon = k\sigma = k$ simpangan baku
untuk suatu bilangan bulat~$k$, berlaku
$$
P(|X - \mu| \geq k\sigma) \leq \frac {\sigma^2}{k^2\sigma^2} = \frac 1{k^2}\ ,
$$
seperti dalam kasus diskret.
\end{example}

\subsection*{Hukum Bilangan Besar}
Dengan Ketaksamaan Chebyshev, kini kita dapat menyatakan dan membuktikan Hukum Bilangan
Besar untuk kasus kontinu.

\begin{theorem}{\bf (Hukum Bilangan Besar)}\index{Hukum Bilangan Besar}
Misalkan $X_1$,~$X_2$, \dots,~$X_n$ membentuk proses percobaan saling bebas dengan
fungsi densitas kontinu~$f$, nilai harapan berhingga~$\mu$, dan
varians berhingga~$\sigma^2$.  Misalkan $S_n = X_1 + X_2 +\cdots+ X_n$ adalah jumlah
$X_i$.  Maka untuk sebarang bilangan real $\epsilon > 0$ berlaku
$$
\lim_{n \to \infty} P\left( \left| \frac {S_n}n - \mu \right| \geq \epsilon
\right) = 0\ ,
$$
atau secara ekuivalen,
$$
\lim_{n \to \infty} P\left( \left| \frac {S_n}n - \mu \right| < \epsilon
\right) = 1\ .
$$
\end{theorem}

Perhatikan bahwa teorema ini belum tentu benar jika $\sigma^2$ tak hingga
(lihat Contoh~\ref{exam 8.2.5}).

Seperti dalam kasus diskret, Hukum Bilangan Besar menyatakan bahwa nilai rata-rata
dari $n$ percobaan saling bebas menuju nilai harapan ketika $n \to \infty$, dalam
arti tepat bahwa, untuk $\epsilon > 0$ yang diberikan, peluang bahwa nilai rata-rata
dan nilai harapan berselisih lebih dari $\epsilon$ menuju~0 ketika $n
\to \infty$.

Sekali lagi, kita tidak menyertakan pembuktiannya karena identik dengan pembuktian dalam kasus diskret.
\subsection*{Kasus Seragam}
\begin{example}
Misalkan kita memilih secara acak $n$ bilangan dari interval $[0,1]$ dengan
distribusi seragam.  Jika $X_i$ menyatakan pilihan ke-$i$, maka
\begin{eqnarray*}
          \mu & = & E(X_i) = \int_0^1 x\, dx = \frac 12\ , \\
     \sigma^2 & = & V(X_i) = \int_0^1 x^2\, dx - \mu^2 \\
              & = & \frac 13 - \frac 14 = \frac 1{12}\ .
\end{eqnarray*}
Jadi,
\begin{eqnarray*}
E \left( \frac {S_n}n \right)  & = & \frac 12\ , \\
V \left( \frac {S_n}n \right)  & = & \frac 1{12n}\ ,
\end{eqnarray*}
dan untuk sebarang $\epsilon > 0$,
$$
P \left( \left| \frac {S_n}n - \frac 12 \right| \geq \epsilon \right) \leq \frac
1{12n \epsilon^2}\ .
$$

Pernyataan ini berarti bahwa jika kita memilih $n$ bilangan secara acak dari $[0,1]$, maka,
dengan peluang lebih besar daripada $1 - 1/(12n\epsilon^2)$, selisih $|S_n/n -
1/2|$ kurang dari~$\epsilon$.  Perhatikan bahwa $\epsilon$ berperan sebagai
besar galat yang bersedia kita toleransi: Jika kita memilih $\epsilon = 0.1$, misalnya,
maka peluang bahwa $|S_n/n - 1/2|$ kurang dari~0.1 lebih besar daripada $1 -
100/(12n)$.  Untuk $n = 100$, nilainya sekitar .92, tetapi jika $n = 1000$, nilainya lebih besar
daripada .99, dan jika $n = 10{,}000$, nilainya lebih besar daripada .999.
\putfig{5.0truein}{PSfig8-2}{Ilustrasi Hukum Bilangan Besar --- kasus seragam.}{fig 8.2}

Kita dapat menggambarkan secara grafis pernyataan Hukum Bilangan Besar untuk contoh ini.
Densitas bagi $A_n = S_n/n$ ditentukan oleh
$$
f_{A_n}(x) = nf_{S_n}(nx)\ .
$$

Kita telah melihat dalam Bagian~\ref{sec 7.2} bahwa kita dapat menghitung densitas
$f_{S_n}(x)$ bagi jumlah $n$ peubah acak seragam.  Dalam Gambar~\ref{fig 8.2}, kita telah
menggunakan hasil ini untuk menggambar densitas $A_n$ bagi berbagai nilai~$n$.  Kita telah
mengarsir daerah tempat $A_n$ terletak di antara .45~dan~.55.  Kita melihat bahwa ketika
$n$ diperbesar, bagian luas total di dalam daerah yang diarsir semakin
besar.  Hukum Bilangan Besar menyatakan bahwa kita dapat memperoleh sebanyak apa pun
bagian luas total di dalam daerah yang diarsir dengan memilih $n$ cukup besar
(lihat juga Gambar~\ref{fig 8.1}).
\end{example}

\subsection*{Kasus Normal}
\begin{example}
Misalkan kita memilih secara acak $n$ bilangan real
menggunakan distribusi normal dengan rata-rata~0 dan varians~1.  Maka
\begin{eqnarray*}
         \mu &=& E(X_i) = 0\ , \\
    \sigma^2 &=& V(X_i) = 1\ .
\end{eqnarray*}
Jadi,
\begin{eqnarray*}
E \left( \frac {S_n}n \right) &=& 0\ , \\
V \left( \frac {S_n}n \right) &=& \frac 1n\ ,
\end{eqnarray*}
dan, untuk sebarang $\epsilon > 0$,
$$
P\left( \left| \frac {S_n}n - 0 \right| \geq \epsilon \right) \leq \frac
1{n\epsilon^2}\ .
$$
Dalam kasus ini, kita dapat membandingkan taksiran Chebyshev bagi $P(|S_n/n -
\mu| \geq \epsilon)$ dalam Hukum Bilangan Besar dengan nilai tepat, karena kita
mengetahui secara tepat fungsi densitas bagi $S_n/n$ (lihat Contoh~\ref{exam 7.12}). 
Perbandingan tersebut ditampilkan dalam Tabel~\ref{table 8.1}, untuk $\epsilon = .1$.  Data
dalam tabel ini dihasilkan oleh program {\bf LawContinuous}.\index{LawContinuous (program)}
\begin{table}
\centering
\begin{tabular}{r|r|r}
$n$ & $P(|S_n/n| \ge .1)$ & Chebyshev \\
\hline
100 & .31731 & 1.00000 \\
200 & .15730 & .50000 \\
300 & .08326 & .33333 \\
400 & .04550 & .25000 \\
500 & .02535 & .20000 \\
600 & .01431 & .16667 \\
700 & .00815 & .14286 \\
800 & .00468 & .12500 \\
900 & .00270 & .11111 \\
1000 & .00157 & .10000 \\
\hline
\end{tabular}
\caption{Taksiran Chebyshev.}
\label{table 8.1}
\end{table}
Di sini kita melihat bahwa taksiran Chebyshev pada umumnya \emx {tidak} begitu 
akurat.
\end{example}

\subsection*{Metode Monte Carlo}
Berikut adalah contoh yang agak lebih menarik.

\begin{example}
Misalkan $g(x)$ adalah fungsi kontinu yang didefinisikan untuk $x \in [0,1]$ dengan nilai dalam
$[0,1]$.  Dalam Bagian~\ref{sec 2.1}, kita menunjukkan cara menaksir luas
daerah di bawah grafik $g(x)$ dengan metode Monte Carlo, yaitu dengan
memilih banyak nilai acak $x$~dan~$y$ yang berdistribusi
seragam, lalu melihat berapa proporsi titik $P(x,y)$ yang jatuh di dalam
daerah di bawah grafik (lihat Contoh~\ref{exam 2.1.2}).
\par
Berikut adalah cara yang lebih baik untuk menaksir luas yang sama (lihat Gambar~\ref{fig 8.3}).  Mari kita pilih
banyak nilai $X_n$ yang saling bebas dan acak dari $[0,1]$ dengan
densitas seragam, tetapkan $Y_n = g(X_n)$, lalu tentukan nilai rata-rata $Y_n$.  Maka
rata-rata ini merupakan taksiran kita bagi luas tersebut.  Untuk melihatnya, perhatikan bahwa jika
fungsi densitas bagi~$X_n$ seragam,
\begin{eqnarray*}
\mu & = & E(Y_n) = \int_0^1 g(x) f(x)\, dx \\
    & = & \int_0^1 g(x)\, dx \\
    & = & \mbox {nilai\ rata-rata\ dari\ }  g(x)\ ,
\end{eqnarray*}
sedangkan variansnya adalah
$$
\sigma^2 = E((Y_n - \mu)^2) = \int_0^1 (g(x) - \mu)^2\, dx < 1\ ,
$$
karena untuk semua $x$ dalam $[0, 1]$, $g(x)$ berada dalam $[0, 1]$, sehingga $\mu$ berada dalam $[0, 1]$, dan
karena itu $|g(x) - \mu| \le 1$.  Sekarang misalkan $A_n = (1/n)(Y_1 + Y_2 +\cdots+ Y_n)$.  Berdasarkan Ketaksamaan
Chebyshev, kita memperoleh
$$
P(|A_n - \mu| \geq \epsilon) \leq \frac {\sigma^2}{n\epsilon^2} < \frac
1{n\epsilon^2}\ .
$$

\putfig{3truein}{PSfig8-3}{Masalah luas.}{fig 8.3} 

Pernyataan ini berarti bahwa agar berada dalam jarak $\epsilon$ dari nilai sebenarnya bagi $\mu = \int_0^1
g(x)\, dx$ dengan peluang sekurang-kurangnya $p$, kita harus memilih $n$ sehingga
$1/n\epsilon^2 \leq 1 - p$ (yakni, sehingga $n \geq 1/\epsilon^2(1 - p)$).  Perhatikan
bahwa metode ini memberi tahu kita seberapa besar $n$ harus dipilih untuk memperoleh ketepatan yang diinginkan.
\end{example}

Hukum Bilangan Besar mensyaratkan agar varians $\sigma^2$ dari densitas
yang mendasarinya berhingga: $\sigma^2 < \infty$.  Dalam kasus ketika syarat ini tidak
terpenuhi, Hukum Bilangan Besar juga dapat gagal.  Berikut sebuah contohnya.

\subsection*{Kasus Cauchy}
\begin{example}\label{exam 8.2.5}
Misalkan kita memilih $n$ bilangan dari $(-\infty,+\infty)$ dengan densitas Cauchy
berparameter $a = 1$.  Kita mengetahui bahwa nilai harapan dan
varians densitas Cauchy tidak terdefinisi (lihat Contoh~\ref{exam 6.23}).  Dalam kasus ini,
fungsi densitas bagi
$$
A_n = \frac {S_n}n
$$
diberikan oleh (lihat Contoh~\ref{exam 7.9})
$$
f_{A_n}(x) = \frac 1{\pi(1 + x^2)}\ ,
$$
yakni, \emx {fungsi densitas bagi $A_n$ sama untuk semua $n$.}  Dalam
kasus ini, ketika $n$ meningkat, fungsi densitas sama sekali tidak berubah dan
Hukum Bilangan Besar tidak berlaku.
\end{example}

\exercises
\begin{LJSItem}

\i\label{exer 8.2.1} Misalkan $X$ adalah peubah acak kontinu dengan rata-rata $\mu =
10$ dan varians $\sigma^2 = 100/3$.  Dengan menggunakan Ketaksamaan Chebyshev, tentukan
batas atas bagi peluang-peluang berikut.
\begin{enumerate}
\item $P(|X - 10| \geq 2)$.

\item $P(|X - 10| \geq 5)$.

\item $P(|X - 10| \geq 9)$.

\item $P(|X - 10| \geq 20)$.
\end{enumerate}

\i\label{exer 8.2.2} Misalkan $X$ adalah peubah acak kontinu dengan nilai yang
berdistribusi seragam pada interval $[0,20]$.
\begin{enumerate}
\item Tentukan rata-rata dan varians $X$.

\item Hitung secara tepat $P(|X - 10| \geq 2)$, $P(|X - 10| \geq 5)$, $P(|X - 10| \geq
9)$, dan $P(|X - 10| \geq 20)$.  Bagaimana jawaban Anda dibandingkan dengan jawaban
Latihan~\ref{exer 8.2.1}?  Seberapa baik Ketaksamaan Chebyshev dalam kasus ini?
\end{enumerate}

\i\label{exer 8.2.3} Misalkan $X$ adalah peubah acak dalam Latihan~\ref{exer
8.2.2}.
\begin{enumerate}
\item Hitung fungsi $f(x) = P(|X - 10| \geq x)$.

\item Sekarang gambarkan grafik fungsi $f(x)$ dan, pada sumbu yang sama, gambarkan grafik
fungsi Chebyshev $g(x) = 100/(3x^2)$.  Tunjukkan bahwa $f(x) \leq g(x)$ untuk semua~$x >
0$, tetapi $g(x)$ bukan hampiran yang terlalu baik bagi~$f(x)$.
\end{enumerate}

\i\label{exer 8.2.100} Misalkan $X$ adalah peubah acak kontinu dengan nilai yang berdistribusi
eksponensial pada $[0,\infty)$ dengan parameter $\lambda = 0.1$.
\begin{enumerate}
\item Tentukan rata-rata dan varians $X$.

\item Dengan menggunakan Ketaksamaan Chebyshev, tentukan batas atas bagi
peluang-peluang berikut: $P(|X - 10| \geq 2)$, $P(|X - 10| \geq 5)$, $P(|X - 10| \geq
9)$, dan $P(|X - 10| \geq 20)$.

\item Hitung peluang-peluang ini secara tepat dan bandingkan dengan batas-batas pada
(b).
\end{enumerate}
\i\label{exer 8.2.101} Misalkan $X$ adalah peubah acak kontinu dengan nilai yang berdistribusi
normal pada $(-\infty,+\infty)$ dengan rata-rata $\mu = 0$ dan varians $\sigma^2 = 1$.
\begin{enumerate}
\item Dengan menggunakan Ketaksamaan Chebyshev, tentukan batas atas bagi
peluang-peluang berikut: $P(|X| \geq 1)$, $P(|X| \geq 2)$, dan $P(|X| \geq 3)$.

\item Luas di bawah kurva normal antara $-1$~dan~1 adalah .6827, antara
$-2$~dan~2 adalah .9545, dan antara $-3$~dan~3 adalah .9973 (lihat tabel dalam
Lampiran~A).  Bandingkan batas-batas Anda pada (a) dengan
nilai-nilai tepat ini.  Seberapa baik Ketaksamaan Chebyshev dalam kasus ini?
\end{enumerate}

\i\label{exer 8.2.102} Jika $X$ berdistribusi normal, dengan rata-rata~$\mu$ dan
varians~$\sigma^2$, tentukan batas atas bagi peluang-peluang berikut dengan menggunakan Ketaksamaan
Chebyshev.
\begin{enumerate}
\item $P(|X - \mu| \geq \sigma)$.

\item $P(|X - \mu| \geq 2\sigma)$.

\item $P(|X - \mu| \geq 3\sigma)$.

\item $P(|X - \mu| \geq 4\sigma)$.
\end{enumerate}
\noindent Sekarang tentukan nilai-nilai tepatnya menggunakan program {\bf NormalArea}\index{NormalArea (program)}
atau tabel normal dalam Lampiran~A, lalu bandingkan.

\i\label{exer 8.2.103} Jika $X$ adalah peubah acak dengan rata-rata~$\mu \ne 0$ dan varians~$\sigma^2$,
definisikan \emx {simpangan relatif} $D$ dari $X$ terhadap rata-ratanya dengan
$$
D = \left| \frac {X - \mu}\mu \right|\ .
$$
\begin{enumerate}
\item Tunjukkan bahwa $P(D \geq a) \leq \sigma^2/(\mu^2a^2)$.

\item Jika $X$ adalah peubah acak dalam Latihan~\ref{exer 8.2.1}, tentukan
batas atas bagi $P(D \geq .2)$, $P(D \geq .5)$, $P(D \geq .9)$, dan $P(D \geq
2)$.
\end{enumerate}

\i\label{exer 8.2.104} Misalkan $X$ adalah peubah acak kontinu dan definisikan {\em
versi baku} $X^*$ dari~$X$ dengan:
$$
X^* = \frac {X - \mu}\sigma\ .
$$
\begin{enumerate}
\item Tunjukkan bahwa $P(|X^*| \geq a) \leq 1/a^2$.

\item Jika $X$ adalah peubah acak dalam Latihan~\ref{exer 8.2.1}, tentukan
batas bagi $P(|X^*| \geq 2)$, $P(|X^*| \geq 5)$, dan $P(|X^*| \geq 9)$.
\end{enumerate}

\i\label{exer 8.2.105}
\begin{enumerate} 
\item Misalkan sebuah bilangan $X$ dipilih secara acak dari $[0,20]$ dengan
peluang seragam.  Tentukan batas bawah bagi peluang bahwa $X$ terletak antara
8~dan~12, dengan menggunakan Ketaksamaan Chebyshev.

\item Sekarang misalkan 20 bilangan real dipilih secara saling bebas dari $[0,20]$
dengan peluang seragam.  Tentukan batas bawah bagi peluang bahwa
rata-rata bilangan tersebut terletak antara 8~dan~12.

\item Sekarang misalkan 100 bilangan real dipilih secara saling bebas dari
$[0,20]$.  Tentukan batas bawah bagi peluang bahwa rata-ratanya terletak
antara 8~dan~12.
\end{enumerate}

\i\label{exer 8.2.106} Nilai seorang mahasiswa pada suatu ujian akhir kalkulus merupakan peubah acak
dengan nilai dalam $[0,100]$, rata-rata~70, dan varians~25.
\begin{enumerate}
\item Tentukan batas bawah bagi peluang bahwa nilai mahasiswa tersebut
terletak antara 65~dan~75.

\item Jika 100 mahasiswa mengikuti ujian akhir tersebut, tentukan batas bawah bagi peluang
bahwa rata-rata kelas terletak antara 65~dan~75.
\end{enumerate}

\i\label{exer 8.2.107} Perusahaan bir Pilsdorff mengoperasikan armada truk di sepanjang jalan
100~mil dari Hangtown ke Dry Gulch dan memelihara sebuah bengkel di tengah perjalanan. 
Setiap truk dapat mengalami kerusakan di titik yang berjarak $X$~mil dari Hangtown,
dengan $X$ peubah acak yang berdistribusi seragam pada $[0,100]$.
\begin{enumerate}
\item Tentukan batas bawah bagi peluang $P(|X - 50| \leq 10)$.

\item Misalkan dalam satu pekan yang buruk, 20 truk mengalami kerusakan.  Tentukan batas bawah
bagi peluang $P(|A_{20} - 50| \leq 10)$, dengan $A_{20}$ adalah rata-rata
jarak dari Hangtown ketika kerusakan terjadi.
\end{enumerate}

\i\label{exer 8.2.12} Satu lembar saham biasa perusahaan bir Pilsdorff
memiliki harga $Y_n$ pada hari kerja ke-$n$ dalam tahun tersebut.  Finn mengamati bahwa
perubahan harga $X_n = Y_{n + 1} - Y_n$ tampaknya merupakan peubah acak dengan rata-rata
$\mu = 0$ dan varians $\sigma^2 =1/4$.  Jika $Y_1 = 30$, tentukan batas bawah bagi
peluang-peluang berikut, dengan asumsi bahwa semua $X_n$ saling bebas.
\begin{enumerate}
\item $P(25 \leq Y_2 \leq 35)$.

\item $P(25 \leq Y_{11} \leq 35)$.

\item $P(25 \leq Y_{101} \leq 35)$.
\end{enumerate}

\i\label{exer 8.2.108} Misalkan seratus bilangan $X_1$,~$X_2$, \dots,~$X_{100}$ dipilih
secara acak dan saling bebas dari $[0,20]$.  Misalkan $S = X_1 + X_2 +\cdots+ X_{100}$
adalah jumlahnya, $A = S/100$ rata-ratanya, dan $S^* = (S - 1000)/(10/\sqrt3)$ jumlah
bakunya.  Tentukan batas bawah bagi peluang-peluang
\begin{enumerate}
\item $P(|S - 1000| \leq 100)$.

\item $P(|A - 10| \leq 1)$.

\item $P(|S^*| \leq \sqrt3)$.
\end{enumerate}

\i\label{exer 8.2.14} Misalkan $X$ adalah peubah acak kontinu yang berdistribusi
normal pada $(-\infty,+\infty)$ dengan rata-rata~0 dan varians~1.  Dengan menggunakan
tabel normal dalam Lampiran~A atau program {\bf NormalArea},
tentukan nilai fungsi $f(x) = P(|X| \geq x)$ ketika $x$ meningkat dari
0~hingga~4.0 dengan langkah~.25.  Perhatikan bahwa untuk $x \geq 0$, tabel memberikan
$ NA(0,x) = P(0 \leq X \leq x)$ sehingga $P(|X| \geq x) = 2(.5 - NA(0,x)$.
Gambarkan dengan tangan
grafik~$f(x)$ menggunakan nilai-nilai ini dan grafik fungsi Chebyshev
$g(x) = 1/x^2$, lalu bandingkan (lihat Latihan~\ref{exer 8.2.3}).

\i\label{exer 8.2.109} Ulangi Latihan~\ref{exer 8.2.14}, tetapi kali ini dengan rata-rata~10 dan
varians~3.  Perhatikan bahwa tabel dalam Lampiran~A menyajikan nilai bagi
peubah normal baku.  Tentukan versi baku $X^*$ dari~$X$, tentukan
nilai $f^*(x) = P(|X^*| \geq x)$ seperti dalam Latihan~\ref{exer 8.2.14}, lalu
skalakan ulang nilai-nilai ini bagi $f(x) = P(|X -10| \geq x)$.  Gambarkan dan bandingkan
fungsi ini dengan fungsi Chebyshev $g(x) = 3/x^2$.

\i\label{exer 8.2.110} Misalkan $Z = X/Y$, dengan $X$~dan~$Y$ memiliki densitas normal dengan rata-rata~0 dan
simpangan baku~1.  Dapat ditunjukkan bahwa $Z$ memiliki densitas Cauchy.
\begin{enumerate}
\item Tulis sebuah program untuk mengilustrasikan hasil ini dengan menggambar diagram batang dari
1000 sampel yang diperoleh dengan membentuk rasio dua hasil normal baku. 
Bandingkan diagram batang Anda dengan grafik densitas Cauchy.  Bergantung pada
bahasa komputer yang digunakan, Anda mungkin perlu atau tidak perlu memberi tahu komputer cara menyimulasikan
peubah acak normal.  Sebuah metode untuk melakukannya dijelaskan dalam Bagian~\ref{sec 5.2}.

\item Kita telah melihat bahwa Hukum Bilangan Besar tidak berlaku bagi
densitas Cauchy (lihat Contoh~\ref{exam 8.2.5}).  Simulasikan banyak
percobaan dengan densitas Cauchy dan hitung rata-rata hasilnya.  Apakah
rata-rata tersebut tampak mendekati suatu limit?  Jika demikian, dapatkah Anda menjelaskan mengapa
hal ini mungkin terjadi?
\end{enumerate}

\i\label{exer 8.2.111} Tunjukkan bahwa, jika $X \geq 0$, maka $P(X \geq a) \leq E(X)/a$.

\i\label{exer 8.2.112} (Lamperti\footnote{Komunikasi pribadi.})
\index{LAMPERTI, J.} Misalkan $X$ adalah peubah acak tak-negatif.
Apakah batas atas terbaik yang dapat Anda berikan bagi $P(X \geq a)$ jika diketahui
\begin{enumerate}
\item $E(X) = 20$.

\item $E(X) = 20$ dan $V(X) = 25$.

\item $E(X) = 20$, $V(X) = 25$, dan $X$ simetris terhadap rata-ratanya.
\end{enumerate}
\end{LJSItem}}
